\section{Matriz analítica proposta}
\label{sec:matriz}

O Gráfico~\ref{fig:heatmap} apresenta a coocorrência documental entre os dez critérios de priorização e os dez instrumentos de apoio à decisão mais frequentes. Desempenho operacional e \textit{framework} aparecem conjuntamente em 88 registros. Informação e dados apresentam 71 coocorrências com \textit{framework} e 26 com BIM. Custo e \textit{framework} aparecem em 55 registros, enquanto vida útil e custo do ciclo de vida aparecem conjuntamente em 13.

\begin{figure}[H]
\centering
\includegraphics[width=0.98\linewidth]{figuras/figura13_heatmap_criterios_metodos_nucleo_final_104.png}
\captiongrafico{Coocorrência entre critérios de priorização e instrumentos de apoio à decisão}
\label{fig:heatmap}
\end{figure}

O Gráfico~\ref{fig:matrizanalitica} reorganiza a coocorrência pelas cinco dimensões de sustentabilidade usadas na matriz analítica. A dimensão técnica-operacional apresenta 95 coocorrências com \textit{framework}, 26 com \textit{decision support} e 26 com BIM. A dimensão institucional apresenta 86 coocorrências com \textit{framework}, seguida da ambiental com 77, da econômica com 55 e da social com 53.

\begin{figure}[H]
\centering
\includegraphics[width=0.98\linewidth]{figuras/figura14_matriz_analitica_dimensao_metodo_nucleo_final_104.png}
\captiongrafico{Matriz de coocorrência entre dimensões de sustentabilidade e instrumentos de apoio à decisão}
\label{fig:matrizanalitica}
\end{figure}

A matriz analítica resultante combina cinco grupos de critérios. A dimensão ambiental reúne energia, emissões, água e resíduos. A econômica reúne custo de manutenção e custo do ciclo de vida. A social reúne segurança, conforto e satisfação dos usuários. A técnica-operacional reúne desempenho, condição física, vida útil e manutenibilidade. A institucional reúne governança, conformidade, informação e capacidade de execução. As linhas da matriz podem ser ponderadas conforme os objetivos da instituição e relacionadas a métodos de pontuação, AHP, TOPSIS, ANP, otimização ou outras formas estruturadas de apoio à decisão.

As coocorrências representam codificações simultâneas no mesmo registro; não estimam associação estatística, efeito, causalidade, peso ou importância dos elementos. Essa síntese transforma os padrões documentais de coocorrência em uma estrutura para pesquisas aplicadas. Sua utilização em instituições públicas de ensino superior requer definição operacional dos indicadores, regras de normalização, processo transparente de ponderação e validação com dados de manutenção, custos, condição dos ativos e percepção dos usuários.

